\documentclass[10pt,a4paper]{amsart}
\usepackage[margin=19mm]{geometry}
\usepackage{amsmath,amssymb,mathtools}
\usepackage[scr=rsfs,cal=euler]{mathalfa}
\usepackage{xfrac}
\usepackage[unicode,hidelinks,bookmarks=false]{hyperref}
\newcommand{\ie}{i.e.\ }
\newcommand{\cf}{cf.\ }
\newcommand{\resp}{respectively\ }
\newcommand{\Subsection}{\subsection}
\providecommand{\nobreakdash}{\nobreak\hbox{-}}
\setlength{\emergencystretch}{2em}
\multlinegap0pt
\usepackage{amssymb}
\usepackage{mathtools}
\usepackage{float}
\usepackage[all]{xy}
\usepackage{tikz}
\newtheorem{lemma}{Lemma}
\newtheorem{proposition}{Proposition}
\newcommand\Z{{\mathbb Z}}
\newcommand\bP{{\mathbb P}}
\newcommand\cW{{\mathcal W}}
\newcommand\cZ{{\mathcal Z}}
\newcommand\fD{{\mathfrak D}}
\newcommand\fZ{{\mathfrak Z}}
\title{On the non-abelian Hodge locus I}
\author{Philip Engel, Salim Tayou}
\date{Source: arXiv:2305.00943v2 (CC BY 4.0)}
\begin{document}
\maketitle

\noindent\emph{Source and licence.} On the non-abelian Hodge locus I. Version arXiv:2305.00943v2 (CC BY 4.0), distributed by arXiv under the Creative Commons Attribution 4.0 International licence, http://creativecommons.org/licenses/by/4.0/, as recorded in the arXiv licence metadata for this exact version and shown on the abstract page https://arxiv.org/abs/2305.00943v2. Full licence notice: \texttt{licenses/arxiv-2305.00943-CC-BY-4.0.txt}.\par\smallskip

\noindent\textit{Editorial scope.} Counted unit: the article's proposition pi1 (source0.tex lines 2118--2121). The article's complete original proof of it is delivered with it (source0.tex lines 2123--2175, one \texttt{proof} environment of 53 lines). No mathematics is written, repaired or re-derived: the statement and the proof are the article's own lines, in the article's own notation, and no retained line is substituted or re-flowed.  Retained uncounted as the source-local context the proof needs: lines 2103--2106.  Nothing was omitted from any retained span, so the package carries no omission record.  No retained line contains a citation, so the wrapper carries no bibliography and no bibliography entry.  No retained line contains a figure environment, an image-inclusion command or drawing code, so no image is delivered and the package declares no asset dependency.  Editorial formatting only, in addition: an AMS \texttt{amsart} preamble that restates the article's own declarations and macros used by the retained lines, namely the environments \texttt{lemma}, \texttt{proposition} on separate counters, together with the standard packages those lines need.  The retained environments print in this document as Lemma 1, Proposition 1 in order of appearance, whereas the article prints the same items with the numbering of its own full document.  The complete native source of the article as distributed in the arXiv e-print for this exact version is bundled unchanged as \texttt{sources/140846/source0.tex}.
\begin{lemma}\label{normal2} Let $Z^\nu\to Z^{\rm red}$ be the
normalization. Then, $\Gamma_Z\subset \Gamma$ is the image
of $\pi_1(Z^\nu)$. It is also the image of $\pi_1(U)$
for any dense open subset $U\subset (Z^{\rm red})_{\rm sm}$. \end{lemma}

\begin{proposition}\label{pi1} Let $Z\in \fD^{\Xi}$ be irreducible
and reduced. The group $\Gamma_Z$ only jumps in size,
in an open neighborhood of $Z\in \fD^{\Xi}$.
\end{proposition}

\begin{proof}
Let $(C,0)\to \fD^{\Xi}$
be an analytic arc, and consider the pullback
family $\fZ\to (C,0)$, with $\fZ_0=Z$. Let $\cW = \fZ^\nu$
be the normalization of the total space. The general fiber $\cW_t$ is normal, so
$\Gamma_{Z_t} = {\rm im}(\pi_1(\cW_t))$ by Lemma
\ref{normal2}.
This is the same group for all $t\in C\setminus 0$
if we assume (as we may) that $\cW$ is a fiber bundle
over $C\setminus 0$. There is a deformation-retraction
$\cW\to \cW_0$ to the central
fiber. Tracing an element of $\pi_1(\cW_t)$ through the
retraction, we get a free homotopy from any $\gamma_t\in \pi_1(\cW_t)$
to an element $\gamma_0\in \pi_1(\cW_0)$.

Conversely, we can lift any element of $\pi_1(\cW_0)$ 
to an element of $\pi_1(\cW_t)$: We have $\pi_1(\cW_0)=\pi_1(\cW) = 
\pi_1(\cW\setminus ((\cW_0)_{\rm sing}\cup \cW_{\rm sing}))$ because $\cW$ is
normal and $(\cW_0)_{\rm sing}\cup \cW_{\rm sing}$ 
has codimension $2$. Thus, any element of $\pi_1(\cW_0)=\pi_1(\cW)$ can
be represented by a loop in $\cW$ avoiding both $(\cW_0)_{\rm sing}$
and $\cW_{\rm sing}$. Then, 
this loop can be deformed off its intersection with $(\cW_0)_{\rm sm}$
as $(\cW_0)_{\rm sm}$ is a locally smooth divisor in $\cW_{\rm sm}$.
So we can represent the loop in $\cW\setminus \cW_0$.
Finally, $\pi_1(\cW\setminus \cW_0)$ is a $\Z$-extension of $\pi_1(\cW_t)$
because it is a fiber bundle over the punctured disk $C\setminus 0$.

Thus, $\Gamma_{Z_t}={\rm im}(\pi_1(\cW_0))$.
%Assume first that $Z$ is generically reduced.
Then the natural morphism
$\cW_0\to \fZ_0=Z$
is a finite birational
morphism because $Z$ is reduced.
Thus, it factors the normalization
$Z^\nu \to \cW_0 \to Z$
and so ${\rm im}(\pi_1(Z^\nu))=
\Gamma_{Z}\subset \Gamma_{Z_t}= {\rm im}(\pi_1(\cW_0)).$ Thus
$\Gamma_Z$ only jumps in size.
%(gets larger as a subset of $\Gamma$)
%in a neighborhood of $Z\in \fD^{\Xi}$. The first statement follows.
%Now, consider $Z$ generically non-reduced. Then
%$\cW_0\to Z^{\rm red}$ is finite, but may
%have degree larger than one---for instance,
%$\cZ$ might be a family of hyperelliptic curves
%collapsing to a doubled $\bP^1$ over $0\in C$,
%while $\cW$ is a smooth family of curves. But
%$\cW_0\to Z^{\rm red}$ is \'etale over
%a dense open subset $U\subset (Z^{\rm red})_{\rm sm}$.
%Thus a finite index subgroup of $\pi_1(U)$ lifts to %$\pi_1(\cW_0)$
%and so $\Gamma_{Z_t}$ contains a finite index subgroup of %$\Gamma_Z$.
%The second statement follows. 
\end{proof}

\end{document}
